\documentclass[a4paper,10pt]{article}

\usepackage{packages/base}
\usepackage{packages/math}
\usepackage{packages/code}
\usepackage{packages/macros}
% ======================================================
% \title{The Title}
% \author{Rustem Sirazetdinov}
% \date{\today}
% ======================================================
\fancyfoot[R]{\thepage}

\hypersetup{
    pdftitle={cv-rustem-sirazetdinov},
    colorlinks=true, % color links?
    allcolors=blue,  % or use | linkcolor | filecolor | urlcolor | ...
}

\usepackage{wrapfig}
\usepackage{changepage}

\newcounter{questioncounter}

\newenvironment{xquestion}{%
    \par\vspace{2.5mm}%
    \refstepcounter{questioncounter}
    \begin{adjustwidth}{0cm}{5cm}%
        \bfseries
        (\thequestioncounter)\
        \ignorespaces
}{%
    \end{adjustwidth}%
    \par\vspace{2.5mm}%
    \ignorespacesafterend
}

\newenvironment{xanswer}{%
    \par\vspace{2.5mm}%
    \begin{adjustwidth}{5cm}{0cm}%
        \ignorespaces
}{%
    \end{adjustwidth}%
    \par\vspace{2.5mm}%
    \ignorespacesafterend
}

\setlength{\parindent}{0cm}
% ======================================================
\begin{document}

\noindent
\begin{minipage}[h]{0.6\textwidth}
    {\bfseries\LARGE Сиразетдинов Рустем Дамирович}\\[5pt]
    РТУ МИРЭА, Прикладная математика и информатика, 2 курс\\
    \url{https://github.com/va9abond}\\
    avesirazetdinov@gmail.com\\
    +7\,(917)\,774--59--97
\end{minipage}
\hfill
\begin{minipage}[h]{0.3\textwidth}
    \hfill\includegraphics[width=3.5cm,height=4.5cm]{face}
\end{minipage}



\begin{xquestion}
    Расскажите, пожалуйста, где и на каком направлении вы учитесь?
    На каком вы сейчас курсе, год выпуска и какой у вас средний балл?
\end{xquestion}

\begin{xanswer}
    Меня зовут Сиразетдинов Рустем. Я студент-математик 2 курса направления
    прикладная математика и информатика (профиль - математическое
    моделирование и вычислительная математика). Год выпуска - 2028.
    Мой средний балл по профильным предметам близок к 5.
    Почти по всем предметам я имею оценку {\itshape отлично} ,
    лишь по дискретной математике (семестр 2)
    я получил оценку {\itshape хорошо}.
    Я вкладываю огромное количество усилий в свое образование,
    поэтому, без ложной скромности,
    могу считать себя одним из лучших студентов своего потока.
\end{xanswer}

\begin{xquestion}
    Какими языками программирования и математическими пакетами
    (Python, C++, MATLAB, Mathematica, Maple, Julia, R) вы владеете?
    Что из этого вы использовали в учебных или собственных проектах?
\end{xquestion}

\begin{xanswer}
    Я наиболее силен в языках C, C++ и Julia.
    Проекты на C
    (\href{https://github.com/va9abond/c-codebase}{ссылка}),
    проекты на C++
    (\href{https://github.com/va9abond/cpp-codebase}{ссылка}),

    \vspace{3mm}
    Я считаю, что любой язык программирования - это инструмент,
    часто не сложный в освоении, он должен выбирается как
    наиболее выразительный для каждой отдельной задачи,
    поэтому у меня нет особых предпочтений в них,
    я с радостью выучу любой из них, если это поможет мне
    в реализации идей.
\end{xanswer}

\begin{xquestion}
    Какие разделы математики и физики вы знаете лучше всего и применяли в реальных
    задачах (например, численные методы, дифференциальные уравнения, линейная
    алгебра, матфизика, механика)?
\end{xquestion}

\begin{xanswer}
    Я считаю математику своей сильной стороной.
    Пока другие студенты сбегают в
    программирование из-за сложности математики, я остаюсь предан ей.
    Возможно, по этой же причене в моем университете редко встречаются
    преподаватели предлагающие проекты связанные с математикой.
    Лично у меня есть лишь один такой проект. Его тема:
    ускорение сходимости последовательностей.
    Моя задача заключалась в том, чтобы представить обзор
    алгоритма экстраполяции Левина, т.е. основываясь на исходных статьях
    Левина, Зиди, Хомейера, я должен был подготовить подробный отчет по
    алгоритму экстраполяции и его вариантам. Сложность заключалась в том, чтобы
    восстановить пропущенные в исходных статьях шаги доказательств и облегчить
    язык статей для их лучшего понимания.

    \vspace{3mm}
    Своей слабой стороной я считаю низкий уровень знания физики. Среди предметов в
    университете физики пока еще не было, ожидается механика и электродинамика
    на 3 курсе. В свою же очередь, я всегда любил физику и участвовал в
    олимпиадах по физике в школе. Я готов вытянуть любой из ее разделов на
    нужный уровень.
\end{xanswer}

\begin{xquestion}
    Помимо этого математического проекта, есть ли у вас другие учебные или личные
    проекты (например, численное решение уравнений, моделирование, обработка
    данных, написание библиотек на C++/Julia)?
    Если да — кратко опишите вашу роль и результат.
\end{xquestion}

\begin{xanswer}
    Моей гордостью среди всего, что я писал на C++, я считаю
    реализацию 
    map
    (\href{https://github.com/va9abond/ac-oop-120-210/blob/5f8d870432ea67c01eeb23eb4ae490cf1b057921/p4-tree/tree.cpp#L204}{ссылка}). ,
    linked\_list в стиле MSVC
    (\href{https://github.com/va9abond/ac-oop-120-210/blob/5f8d870432ea67c01eeb23eb4ae490cf1b057921/p4-tree/tree.cpp#L778}{ссылка}).
    Эти проекты нельзя назвать серьезными,
    но я вложил в них много усилий и получил много ценного опыта
    во время их написания.
    Роль - разработчик (единственный, т.к. это
    мои личные проекты). Итог - хорошие учебные копии контейнеров.

    \vspace{3mm}
    Дополнительно я обучаюсь в Школе 21 на направлении
    {\itshape Разработка/Backend}.
\end{xanswer}

\begin{xquestion}
    Работаете ли вы в Linux (какие дистрибутивы),
    используете ли Git, системы сборки (CMake/Make),
    отладчики (gdb, lldb) или профилировщики?
\end{xquestion}

\begin{xanswer}
    Я использую GNU/Linux как свою основную операционную систему, а конкретно
    EndeavourOS. У меня нет никаких проблем с
    использованием любого из дистрибутивов GNU/Linux.
    Git и Make мои добрые друзья в любом из проектов.
    Как раз в данный момент учусь пользоваться отладчиком gdb.
\end{xanswer}

\begin{xquestion}
    Были ли у вас научные доклады, участие в конференциях
    (студенческих или научных), публикации или олимпиады
    по математике/программированию? Даже
    если без дипломов — сам опыт выступления важен.
\end{xquestion}

\begin{xanswer}
    У меня нет опыта участия в конференциях и нет никаких
    публикаций. Ранее я участвовал в соревнованиях по
    спортивному программированию, но это не более чем просто хобби.
\end{xanswer}

\begin{xquestion}
    Какой у вас уровень английского языка
    (чтение научных статей, переписка, устная речь)?
    Сдавали ли какой-либо экзамен (TOEFL, IELTS)?.
\end{xquestion}

\begin{xanswer}
    Я свободно читаю документацию или технические сайты на английском языке. Чтение
    научных статей дается сложнее - без переводчика бывает не обойтись,
    из-за обилия терминов, которые не всегда совпадают с привычными русскоязычными вариантами. Могу
    вести переписку и устный разговор на несложные темы.
    Экзамены не сдавал.
\end{xanswer}

\begin{xquestion}
    Расскажите о каком-нибудь случае, когда вы столкнулись с трудной математической
    или алгоритмической задачей (учебной, олимпиадной или из проекта) и
    самостоятельно нашли способ её решить. Что это была за задача и каков был
    ваш подход?
\end{xquestion}

\begin{xanswer}
    Например, совсем недавно мне встретилась в домашней работе про дифференциальным
    уравнениям задача. Мне нужно было нарисовать траектории автономной системы
    на фазовой плоскости. Я несколько раз штурмовал задачу,
    и в процессе решения ввел для себя понятие
    {\itshape ось(кривая) экстремумов} - эта прямая (или кривая),
    на которой лежат экстремумы фазовых траекторий. Тогда оно оказалось
    важным и помогло мне правильно проанализировать и изобразить вид фазовых траекторий.

    \vspace{3mm}
    Другой пример, из недавнего проекта Школы 21. Мне нужно было реализовать
    структуру матрицы и базовую алгебру над ними. По моему скромному опыту,
    самая сложная задача в разработке программного обеспечения - это
    продумывание его архитектуры (API design): то, как компоненты программного
    обеспечения будут взаимодействовать между собой, как сделать максимально
    простые и гибкие пайплайны для общения с внешними библиотеками, причем так
    чтобы все это было покрыто простым и понятным интерфейсом.
    Предложенный авторами интерфейс мне совершенно не нравился
    из-за отсутствия гибкости. Еще одна проблема была связана с
    огромным количеством проверок входных данных, которые, как и любое
    ветвление, замедляли код. Путем проб и ошибок, долгих обсуждений вариантов с ИИ я
    пришел, как мне кажется, к хорошему решению.
\end{xanswer}

\begin{xquestion}
    Какое лабораторное или специализированное ПО вы использовали (например, Wolfram
    Mathematica, Maple, MATLAB/Simulink, COMSOL, Ansys, LabVIEW, Jupyter,
    системы компьютерной вёрстки (LaTeX) или другое)? Для техника/стажёра в НИИ
    это бывает важно.
\end{xquestion}

\begin{xanswer}
    Был опыт работы работы с блокнотами Jupyter. Wolfram|Alpha (но не Wolfram
    Mathematica) - мой хорший друг. Отлично владею версткой в системе Latex,
    которая уже давно стала мне гораздо ближе чем MS Word или Google Docs.
\end{xanswer}

\begin{xquestion}
    Как вы организуете свою работу над сложной задачей, которая не решается за один
    день? Например, как планируете время, фиксируете промежуточные результаты,
    работаете с ошибками?
\end{xquestion}

\begin{xanswer}
    Если я встречаю действительно сложную задачу, то решаю ее так: в первом подходе
    я могу мусолить задачу довольно долго, пытаясь взять ее силой (техникой).
    Если задача не решилась, то я откладываю ее. Возвращаясь к задаче снова, я
    уже пытаюсь взять ее идейно (т.е. подключаю все свою лень и смекалку, чтобы
    придумать решение не в лоб). Суть в том, чтобы вернуться к задаче несколько
    раз, каждый раз пытаясь пройти по новому пути. Стараюсь не сидеть за
    задачей до посинения, хотя случается и такое. Если задача не берется со 2-3
    идейной попытки, то я обращаюсь к преподавателю или к всезнающему (т.е. к ИИ).
\end{xanswer}

\begin{xquestion}
    Есть ли у вас опыт работы в команде над научным или инженерным проектом
    (учебным, в Школе 21 или ещё где-то)? Если да, то какую роль вы обычно
    брали на себя?
\end{xquestion}

\begin{xanswer}
    Я работал над несколькими командными проектами в Школе 21.
    Был тимлидом и главным разработчиком.
    Роль тимлида включала в себя организацию всего
    процесса разработки и взаимодействия между разработчиками. Я понял, что
    занятие орг. вопросами мне не близко. Гораздо большее мне откликается
    написание кода - непосредственная работа над самим проектом, когда я
    подходить к решению творчески, придумывать нестандартные и
    максимально эффективные решения. Таким образом,
    сейчас для меня наиболее близка роль исполнителя задач.
\end{xanswer}

\begin{xquestion}
    Почему вы хотите работать именно в научно-исследовательском институте
    (а не в коммерческой разработке)?
    И какие у вас амбиции в профессиональном плане?
\end{xquestion}

\begin{xanswer}
    Меня всегда поражали своей гениальностью машины, самолеты, подводные и
    надводные корабли, Большой Адронный Коллайдер, ракеты и компьютеры - те
    вещи, которые состоят из бесчисленного количества деталей и приборов,
    работающих друг с другом слаженно и вместе составляющих сложные
    удивительные системы. И со школьных времен моей мечтой было - стать
    выдающимся математиком, чтобы внести свой вклад в создание таких же сложных
    и полезных для человечества механизмов.

    \vspace{3mm}
    Через 5-6 лет я рассчитываю стать опытным и уверенным специалистом.
    Через 7-10 лет - ключевым ведущим сотрудников в компании и специалистом
    топ-уровня через 14-19 лет. И в далеком будущем
    хотел бы создать собственное конструкторское бюро.
\end{xanswer}


\end{document}
